% Chapter 6 — Implementation
% This chapter is self-contained and focuses strictly on implementation details.
\chapter{Implementation}

\section{Development Process}
\subsection{Approach and Stack}
The app is a React (Vite) SPA that calls a Node.js/Express API backed by PostgreSQL. Features were built end\-to\-end in vertical slices (wizard, plans, calendar, logging, Strava). We prioritised stable JSON contracts, small stateful components, and controller/utility isolation.

\textbf{Stack}: React 18, Vite, React Router, \texttt{date\-fns}, Express, PostgreSQL, Python (scikit\-learn KMeans), OpenAI API, Strava OAuth. Configuration is centralised in \verb|client/src/config.js| and environment variables; the SQL schema lives in \verb|server/db/schema.sql|.

\section{Feature Implementations}
\subsection{User Authentication}
\subsubsection{Purpose}
Provide sign-in/registration, session handling, and secure API access. Social login and Strava OAuth are integrated to link an athlete account for activity import.

\subsubsection{Implementation}
Authentication routes are defined in \verb|server/routes/authRoutes.js| and registered in \verb|server/app.js|. The client exposes login/registration pages and a Strava OAuth callback component. For Strava, tokens are stored in \verb|strava_tokens| and refreshed when expired.

\subsubsection{Key Logic}
OAuth callback exchanges codes for tokens; middleware guards routes; the client stores \verb|userId| and a bearer token in \verb|localStorage|.

\subsection{Goal Management}
\subsubsection{Purpose}
Capture the runner’s training goal: distance, target time, preferred training days, and weekly volume.

\subsubsection{Implementation}
A multi\-step wizard in \verb|client/src/pages/TrainingPlanCalendarPage.jsx| captures distance, target time, preferred days and weekly volume. The client persists the goal through \verb|client/src/api/goalsAPI.jsx| to \verb|/api/goals|, which stores records in \verb|goals|. The most recent goal for a user is used as the active goal.

\subsubsection{Key Workflows}
On load the client fetches goals and selects the most recent; the wizard accumulates fields in component state and saves once at confirmation.

\subsection{AI-Powered Training Plan Generation}
\subsubsection{Purpose}
Transform a goal and recent training data into a week-by-week plan with run types, distances, and target paces.

\subsubsection{Controller Flow}
The backend exposes plan endpoints in \verb|server/routes/planRoutes.js|, handled by \verb|server/controllers/planController.js|:
\begin{enumerate}
  \item Fetching the active goal and recent \verb|training_records|.
  \item Calling utilities (\verb|utils/aiPlanGenerator.js|, \verb|utils/kmeans_predict.js|, \verb|utils/cluster_info.js|) that use KMeans cluster priors exported from Python (\verb|ml/export_kmeans_to_json.py|) to shape weekly intensities and run mixes.
  \item Returning a normalized payload with \verb|plan.details| (\verb|week|, \verb|day|, \verb|type|, \verb|distance|/\verb|distance_km|, \verb|target_pace|, optional \verb|note| and \verb|explanation|). When persisted, rows are written to \verb|plan_details|.
\end{enumerate}

\subsubsection{Clustering Pipeline (Offline)}
Historical runs are clustered offline to obtain intensity archetypes that inform plan structure.
\begin{enumerate}
  \item \textbf{Data extraction}: \verb|ml/clustering.py| queries \verb|training_records| for \(\{\)distance\_km, duration\_minutes, pace, average\_cadence, average\_heartrate\(\}\)) and drops rows with missing fields.
  \item \textbf{Model fitting}: scikit\-learn KMeans is trained with \(k=10\) and serialized to \verb|kmeans_model.pkl|. Representative samples per cluster are exported for inspection; the full model is converted to JSON with \verb|ml/export_kmeans_to_json.py|.
  \item \textbf{Serving artifact}: the JSON contains \verb|feature_names|, \verb|centroids|, and metadata (\verb|n_clusters|, \verb|inertia|). The Node utility \verb|server/utils/kmeans_predict.js| loads this artifact and computes nearest\-centroid assignments using Euclidean distance; optional standardization is supported if scaler statistics are included.
\end{enumerate}

\subsubsection{Prompt Construction and GPT Invocation}
Two prompt builders exist in \verb|planController.js|:
\begin{itemize}
  \item \verb|buildSystemPrompt|: used when only baseline metrics are available; includes goal, preferred days, weekly volume, and average easy pace.
  \item \verb|buildSystemPromptWithRecentRuns|: used when prior plan data exists; embeds the last 10 runs (optionally annotated with KMeans cluster metadata from \verb|cluster_info.js|) to constrain GPT's plan with recent performance context.
\end{itemize}
Both prompts instruct the model to return \emph{raw JSON only} with required fields and guardrails (3\+1 microcycle, \(\leq 10\%\) weekly increase, weekend long run). The controller calls \verb|openai.responses.create| (model \verb|gpt-4.1-nano|), saves the raw JSON to \verb|server/plans/| for traceability, parses it, and persists rows to \verb|training_plans| and \verb|plan_details|.

\subsubsection{Data Ingestion and Feature Engineering}
The pipeline that powers clustering and plan personalization begins with structured ingestion of workout records.
\begin{enumerate}
  \item \textbf{Import from external sources}: Strava activities are normalized in \verb|server/controllers/stravaController.js| and stored in \verb|training_records| with numeric fields (distance\_km, duration\_minutes, pace, average\_cadence, average\_heartrate). Heart rate is rounded to integers at insert to keep a consistent type.
  \item \textbf{Local scripts}: \verb|import_training_records.py| supports ad\-hoc loading or transformation of CSV exports into the same schema for development and experimentation.
  \item \textbf{Derived fields}: Pace is consistently represented as minutes per kilometer, enabling simple averages and comparisons across runs and clusters.
\end{enumerate}

\subsubsection{KMeans Model Training (\verb|ml/kmeans.py| and \verb|ml/clustering.py|)}
The KMeans step discovers archetypal effort profiles from historical runs.
\begin{enumerate}
  \item \textbf{Query}: A SQL selection over \verb|training_records| retrieves the five features used by the model: distance\_km, duration\_minutes, pace, average\_cadence, average\_heartrate.
  \item \textbf{Cleaning}: Rows with any missing values are dropped to ensure a well\-conditioned fit.
  \item \textbf{Training}: scikit\-learn \verb|KMeans(n_clusters=10)| is fitted with a fixed \verb|random_state| for reproducibility. The trained estimator is persisted as \verb|kmeans_model.pkl|.
  \item \textbf{Inspection}: Clustered data and uniform samples per cluster are written out as CSV to manually inspect cluster separation and confirm face validity.
\end{enumerate}

\subsubsection{Model Export and Online Prediction}
For use inside Node.js, the Python model is exported to JSON by \verb|ml/export_kmeans_to_json.py|:
\begin{itemize}
  \item The artifact encodes \verb|feature_names| and \verb|centroids| (optionally scaler statistics), plus metadata (\verb|n_clusters|, \verb|inertia|).
  \item \verb|server/utils/kmeans_predict.js| loads this JSON once and exposes \verb|predict(sample)|, which orders features according to \verb|feature_names|, optionally standardizes them, and returns the nearest centroid index via Euclidean distance.
\end{itemize}
During plan generation for returning users, the controller annotates the last ten runs with cluster indices and human\-readable descriptions from \verb|server/utils/cluster_info.js|. These annotations are inserted verbatim in the prompt used to bound GPT outputs.

\subsubsection{Plan Persistence and Retrieval}
When a plan is generated, the raw JSON is stored under \verb|server/plans/| as an audit trail. In the same transaction, a new row is inserted into \verb|training_plans| (start and end dates inferred for 16 weeks) and each JSON entry is inserted into \verb|plan_details| with schema\-conformant fields. The \verb|get-plan| endpoint reconstructs a plan for the client by joining \verb|training_plans| with \verb|plan_details| and returning \verb|details| as an ordered array.

\paragraph{Key Algorithms and Logic}
\begin{itemize}
  \item \textbf{Cadence–pace–distance heuristics}: Target paces are bounded by recent performance features.
  \item \textbf{3+1 microcycle}: Three build weeks followed by a deload week, enforced when generating week sequences.
  \item \textbf{Progression safeguards}: Weekly total distance constrained (\(\leq 10\%\) increase) and run types alternated (easy/interval/long/rest).
\end{itemize}

\subsection{Training Plan Calendar and Weekly View}
\subsubsection{Purpose}
Present the plan by week, allow navigation across weeks, and display per-run details with user-entered performance.

\subsubsection{Implementation}
The weekly view is implemented directly in \verb|TrainingPlanCalendarPage.jsx| using a flex grid (one column per weekday) and \verb|date-fns| for week boundaries. Each cell renders type, distance and target pace; details expand to show notes and explanations. A fixed overlay with a centered GIF is displayed while generating or updating plans.

\subsubsection{Training Wizard (Goal Capture)}
The wizard embedded in \verb|TrainingPlanCalendarPage.jsx| advances through discrete steps stored in local component state (\verb|step|), accumulating inputs in a single \verb|data| object:
\begin{enumerate}
  \item \textbf{Category \/ Distance}: Users select the motivation (race or distance) and the target distance (5K, 10K, Half, Marathon, Ultra).
  \item \textbf{Target Time}: A time string in \verb|HH:MM:SS| or \verb|MM:SS| format is captured.
  \item \textbf{Preferred Days}: A multi\-select over weekdays builds the \verb|days| array.
  \item \textbf{Weekly Distance}: A numeric target in kilometers.
  \item \textbf{Confirmation}: A review screen summarizes selections and triggers goal persistence via \verb|saveGoal|; on success it calls the plan generation endpoint.
\end{enumerate}
All steps are pure client state transitions; only the final confirmation performs network IO.

\subsubsection{Key Interactions}
\begin{itemize}
  \item Week navigation adjusts \verb|week| state; visible runs are filtered as \verb|plan.filter(r => r.week === week)|.
  \item A modal offers run details, while a side section exposes actions for adding a new goal or updating a plan.
\end{itemize}

\subsection{Performance Logging}
\subsubsection{Purpose}
Record actual workout results (distance, duration, pace, cadence, heart rate) for plan adherence and feedback.

\subsubsection{Implementation}
Client helpers in \verb|client/src/api/performanceAPI.js| call \verb|POST /api/plans/save-performance| and \verb|GET /api/plans/get-performance|. Inputs are bound per weekday; save triggers parallel POSTs only for visible runs. On load, existing values prefill fields. Pace is computed client\-side from distance and duration and rendered read\-only as \verb|min:sec/km|.

\subsubsection{Week\/Day to Date Resolution}
For persistence, the server resolves a specific calendar date from (user, week, day): the latest \verb|training_plans| row provides \verb|start_date|; \verb|run_date| is computed as \verb|start_date + 7*(week-1) + dayIndex|. Upserts into \verb|training_records| are keyed by \verb|user_id|, \verb|run_date|, and \verb|record_type='run'|, ensuring idempotent saves.

\paragraph{Data Model}
The \verb|training_records| table stores granular workout data with indices on \verb|user_id| and \verb|run_date| for retrieval by day and week.

\subsection{Strava Integration}
\subsubsection{Purpose}
Import historical and ongoing activities for a richer training history.

\subsubsection{Implementation}
\verb|server/controllers/stravaController.js| manages OAuth token storage (\verb|strava_tokens|), refresh, and activity ingestion into \verb|training_records|. Heart rate values are cast to integers before insert. The dashboard uses endpoints in \verb|planRoutes| to fetch aggregated weekly stats (\verb|/api/plans/dashboard-stats|), recent activities (\verb|/api/plans/recent-activities|), and the demo Strava id (\verb|/api/plans/strava-user-id|).

\subsection{Dashboard}
\subsubsection{Purpose}
Provide a concise overview of goals, weekly progress, recent activities, and key performance indicators.

\subsubsection{Implementation}
\verb|client/src/api/dashboardAPI.js| centralizes calls for the Strava id, weekly stats, and recent activities using a shared \verb|API_URL|. \verb|DashboardPage.jsx| consumes these helpers and renders progress and activity lists with minimal client\-side aggregation.

\section{Technologies and Tools}
\subsection{Stack Summary}
\begin{itemize}
  \item \textbf{Frontend}: React, Vite, React Router, \texttt{react-big-calendar}, \texttt{date-fns}.
  \item \textbf{Backend}: Node.js (Express), modular routers/controllers, middleware error handling.
  \item \textbf{Database}: PostgreSQL with normalized tables for users, goals, plans, plan details, Strava tokens, and training records.
  \item \textbf{ML Tooling}: Python scikit-learn (KMeans), export to JSON, consumed by Node utilities.
  \item \textbf{3rd-party}: Strava API for OAuth and activity data.
\end{itemize}

\subsection{Rationale for Choices}
React/Vite enable rapid iteration; Express provides an explicit JSON API surface; PostgreSQL gives relational integrity and indexed aggregations; KMeans yields explainable priors without heavy inference; Strava provides OAuth and reliable activity data.

\section{Implementation Challenges}
\subsection{Data Contract Stability}
Endpoints consistently return normalized shapes (e.g., \verb|plan.details| with \verb|week|/\verb|day|/\verb|type|), limiting component conditionals. When fields evolved (e.g., \verb|distance| vs. \verb|distance_km|), mapping was centralized in utilities and controllers rather than spread through UI code.

\subsection{Calendar Semantics}
ISO week alignment and week indexing are derived from a single \verb|weekStartsOn=1| constant via \verb|date-fns| utilities, preventing off\-by\-one inconsistencies.

\subsection{Unit Conversion and Formatting}
Users enter decimal minutes; the UI converts to seconds, derives pace, and renders \verb|min:sec|. Rollover (\(\geq 60\) seconds) is handled before storing to avoid drift.

\subsection{External API Constraints}
Strava rate limits and token expiry are mitigated by caching tokens in \verb|strava_tokens| and refreshing on demand. The dashboard tolerates missing data by falling back to empty aggregates.

\section{Coding Practices}
\subsection{Standards and Conventions}
Identifiers are descriptive (functions as verbs, values as nouns). Conditional logic is extracted into helpers where clarity improves. Comments capture intent, not restatements. Styling is component\-scoped; larger components are split at state boundaries.

\subsection{Maintainability Strategies}
\begin{itemize}
  \item Centralized API wrappers (\verb|client/src/api/*|) to avoid duplicated fetch logic.
  \item Separation of concerns: routing in \verb|routes/|, orchestration in \verb|controllers/|, computation in \verb|utils/|.
  \item Data access helpers encapsulate SQL and mapping into controller-level functions.
  \item Defensive defaults for optional fields to keep UI robust against partial data.
\end{itemize}

\section{Key Code Examples}
\subsection{Plan Generation Trigger (Client)}
\begin{lstlisting}[language=JavaScript, caption={Triggering AI plan generation from the wizard}]
const handleWizardComplete = async (answers) => {
  setLoading(true);
  const res = await fetch('/api/plans/generate-ai', {
    method: 'POST',
    headers: { 'Content-Type': 'application/json',
               'Authorization': `Bearer ${localStorage.getItem('token')}` },
    credentials: 'include',
    body: JSON.stringify({ userId, ...answers })
  });
  const data = await res.json();
  setPlan(data.plan.details || []);
  setLoading(false);
};
\end{lstlisting}

\subsection{Offline Clustering (Python)}
\begin{lstlisting}[language=Python, caption={Training KMeans on historical runs}]
import pandas as pd
from sklearn.cluster import KMeans
import psycopg2, pickle

conn = psycopg2.connect(...)
df = pd.read_sql(
    """
    SELECT distance_km, duration_minutes, pace,
           average_cadence, average_heartrate
    FROM training_records
    WHERE record_type='run'
    """,
    conn,
)
df = df.dropna()

kmeans = KMeans(n_clusters=10, random_state=42)
df['cluster'] = kmeans.fit_predict(df)

with open('kmeans_model.pkl', 'wb') as f:
    pickle.dump(kmeans, f)
\end{lstlisting}

\subsection{Export Model to JSON (Python)}
\begin{lstlisting}[language=Python, caption={Converting the fitted model to a JSON artifact}]
import json, pickle
FEATURE_NAMES = [
  'distance_km','duration_minutes','pace','average_cadence','average_heartrate'
]
with open('kmeans_model.pkl','rb') as f:
    model = pickle.load(f)
export = {
  'n_clusters': int(model.n_clusters),
  'inertia': float(model.inertia_),
  'feature_names': FEATURE_NAMES,
  'centroids': model.cluster_centers_.tolist(),
}
open('kmeans_model.json','w').write(json.dumps(export, indent=2))
\end{lstlisting}

\subsection{Online Prediction (Node)}
\begin{lstlisting}[language=JavaScript, caption={Assigning a run to the nearest centroid}]
const { feature_names, centroids } = JSON.parse(fs.readFileSync(modelPath,'utf8'));

function predict(sample) {
  const x = feature_names.map(f => Number(sample[f]));
  let best = { idx: 0, d: Infinity };
  for (let i = 0; i < centroids.length; i++) {
    const c = centroids[i];
    const d = Math.sqrt(c.reduce((s, ci, j) => s + (x[j]-ci)**2, 0));
    if (d < best.d) best = { idx: i, d };
  }
  return best.idx;
}
\end{lstlisting}

\subsection{Prompt with Recent Runs and Clusters (Node)}
\begin{lstlisting}[language=JavaScript, caption={Building the GPT system prompt with clustered runs}]
const runs = await getUserRecentRunsFromTrainingRecords(userId, 10);
const kmeans = loadKMeansModel(kmeansPath);
const annotated = runs.map(r => ({
  ...r,
  cluster: kmeans.predict(r),
  clusterInfo: clusterInfo[kmeans.predict(r)] || null,
}));
const systemPrompt = buildSystemPromptWithRecentRuns({
  username, goalDistance: goal.distance, goalTime: goal.target_time,
  runDays, weeklyDistance, recentRuns: annotated,
});
\end{lstlisting}

\subsection{Weekly Performance Auto-Pace}
\begin{lstlisting}[language=JavaScript, caption={Deriving pace from distance and duration}]
if (field === 'distance_km' || field === 'duration_minutes') {
  const d = parseFloat(updated[day]?.distance_km);
  const t = parseFloat(updated[day]?.duration_minutes);
  if (d > 0 && t > 0) {
    const totalSeconds = Math.round(t * 60);
    const paceSeconds = Math.round(totalSeconds / d);
    const paceMin = Math.floor(paceSeconds / 60);
    const paceSec = paceSeconds % 60;
    updated[day].pace = `${paceMin}:${paceSec.toString().padStart(2, '0')}`;
  } else {
    updated[day].pace = '';
  }
}
\end{lstlisting}

\subsection{Database Schema (Excerpt)}
\begin{lstlisting}[language=SQL, caption={Tables for goals, plans, details, and training records}]
CREATE TABLE goals (
  id BIGSERIAL PRIMARY KEY,
  user_id BIGINT NOT NULL REFERENCES users(id) ON DELETE CASCADE,
  distance VARCHAR(50) NOT NULL,
  target_time VARCHAR(50) NOT NULL,
  created_at TIMESTAMP DEFAULT NOW()
);

CREATE TABLE plan_details (
  id BIGSERIAL PRIMARY KEY,
  training_plan_id BIGINT NOT NULL REFERENCES training_plans(id) ON DELETE CASCADE,
  week INTEGER NOT NULL,
  day VARCHAR(10) NOT NULL,
  type VARCHAR(50) NOT NULL,
  distance DECIMAL(5,2),
  target_pace TEXT,
  note TEXT,
  explanation TEXT
);
\end{lstlisting}

\subsection{Dashboard API Wrapper}
\begin{lstlisting}[language=JavaScript, caption={Consistent API layer using a shared base URL}]
import { API_URL } from '../config';
export const getDashboardStats = async (userId) => {
  const token = localStorage.getItem('token');
  const res = await fetch(`${API_URL}/plans/dashboard-stats?userId=${userId}`, {
    headers: { 'Authorization': `Bearer ${token}`, 'Content-Type': 'application/json' },
    credentials: 'include'
  });
  if (!res.ok) throw new Error('Failed to fetch dashboard stats');
  return { success: true, data: await res.json() };
};
\end{lstlisting}

% Requirements coverage is described in earlier chapters; omitted here per thesis structure.

\section{Summary}
This chapter detailed the concrete implementation of a full-stack running coach application. The solution integrates a React SPA, an Express API, and a PostgreSQL store; augments user goals with ML-informed training logic; and incorporates Strava for activity data. Each feature is realized through cohesive modules: goal management, plan generation with week structuring and pacing rules, calendar visualization, performance logging, and a metrics-focused dashboard. The technology choices and coding practices emphasize clarity, evolvability, and user-centric responsiveness, directly supporting the project’s goals of delivering personalized and actionable training guidance. 
